that it follows from 7.16.2 that @@M0@@ is isomorphic to a quotient of a subgroup of @@M1@@ (@@M2@@ proper and smooth); in particular, its rank is again bounded above by the rank of the latter.

\sgaStableTarget{sga6.e10.para.7-17-4}\textbf{7.17.4.} When @@M3@@ is regular at @@M4@@ and @@M5@@ is maximal, one can construct @@M6@@ by blowing up @@M7@@ in the localization of @@M8@@ at @@M9@@, and localizing at the generic point of the fiber of @@M10@@ in the blow-up. One then obtains for @@M11@@ a \emph{pure extension} of @@M12@@. Now it is not difficult to show that, when @@M13@@ is a pure extension of a field @@M14@@, for every noetherian scheme @@M15@@ over @@M16@@ the canonical homomorphism @@M17@@ is an isomorphism, by representing @@M18@@ as the inductive limit of the coordinate rings of affine open subsets of schemes of the form @@M19@@, and using (IX 1.6). Thus one concludes that (7.17.1) can be regarded as a canonical homomorphism
\sgaStableTarget{sga6.e10.eq.7-17-4-1}@@M20@@
But note that one cannot deduce from this a homomorphism @@M21@@, for when one passes to a finite extension @@M22@@ of @@M23@@ and normalizes @@M24@@ in @@M25@@, the regularity hypothesis on @@M26@@ is lost.

Another way to define a homomorphism (7.17.4.1) is to reduce to the case where @@M27@@ is local with closed point @@M28@@ (which is trivial), and to consider an increasing sequence of \emph{regular} closed subschemes @@M29@@ of @@M30@@, with generic points @@M31@@, such that @@M32@@ (@@M33@@). Since for every @@M34@@, @@M35@@, @@M36@@ is a discrete valuation ring with residue field @@M37@@ and fraction field @@M38@@, there is a specialization homomorphism
@@M39@@
